\chapter{\abstractname}

Tensor network contractions are a fundamental computational primitive in applications such as quantum circuit simulation, many-body physics, computational chemistry, and machine learning. While considerable research has focused on minimizing the arithmetic cost of tensor contractions through optimized contraction ordering, peak memory usage and memory movement remain major bottlenecks on modern computing systems. Existing approaches typically treat tensor permutations as local implementation details, often leading to unnecessary data movement and increased memory consumption throughout the contraction process.

This thesis presents a memory-aware post-contraction path discovery optimization framework that treats tensor layouts as a first-class optimization objective after the contraction path has been determined. A greedy permutation planning algorithm is proposed that identifies memory-critical intermediate tensors, preserves their layouts whenever beneficial, and propagates compatible tensor layouts throughout the contraction tree to reduce redundant permutations. The approach is implemented as a modular Python library integrated with the \ac{TCCG} and the \ac{HPTT} library, enabling efficient execution using optimized native tensor contraction and transposition kernels.

The proposed method is evaluated through both theoretical analysis and practical experiments on thousands of randomly generated tensor networks. The theoretical evaluation demonstrates substantial reductions in predicted peak memory usage, achieving improvements of up to approximately 42\% while maintaining computational efficiency. Practical experiments show more modest gains due to the extensive low-level memory optimizations already present in state-of-the-art tensor contraction frameworks. Nevertheless, the proposed approach performs competitively with existing implementations while providing a modular optimization layer that can be integrated with future tensor contraction systems.

Overall, this work demonstrates that global tensor layout optimization is a promising complement to traditional contraction path optimization and represents an important step toward more memory-efficient tensor network execution.
